\documentclass[11pt]{article}
\usepackage[T1]{fontenc}
\usepackage{lmodern}
\usepackage[margin=1in]{geometry}
\usepackage{amsmath,amssymb,amsthm,mathtools,booktabs,graphicx,microtype}
\usepackage[numbers]{natbib}
\usepackage[colorlinks=true,allcolors=blue]{hyperref}
\newtheorem{theorem}{Theorem}
\newtheorem{proposition}[theorem]{Proposition}
\newtheorem{lemma}[theorem]{Lemma}
\newtheorem{corollary}[theorem]{Corollary}
\theoremstyle{definition}\newtheorem{definition}[theorem]{Definition}
\theoremstyle{remark}\newtheorem{remark}[theorem]{Remark}
\newcommand{\R}{\mathbb R}
\newcommand{\X}{\mathcal X}
\newcommand{\CH}{\mathrm{CH}}
\newcommand{\clamp}{\operatorname{clamp}}
\title{Marginal-Price Support for Fully Replenished Electric-Bus Fleets}
\author{Nathan Cho}
\date{Research draft, 29 October 2026}
\begin{document}
\maketitle
\begin{abstract}
Continuous charging does not make service assignments divisible. We study
whether a complete electric-bus fleet minimizes its operating cost and
electricity bill at the marginal supply prices associated with its own load.
A two-service construction has no supported physical plan despite restoring
every used battery to its initial full inventory. We characterize exactly
when connector capacity and vehicle cost create this obstruction. Additional
nonnegative-energy movements preserve it under the original equipment and
charging windows, whereas an explicitly enlarged battery choice can remove it.
A higher fee for that added battery gives a counterexample to a natural
sufficiency test based on the old convexified value. Replication separates
vanishing planning gaps from a whole owner's absolute deviation incentive;
participants with reserved individual resources have vanishing incentives.
The supporting-price machinery is classical. The contribution developed here
is the physical construction and its conditional stability and erasure
characterizations. A public 37-service case has informative exact bounds but
its optimal-fleet price support remains unresolved. These analytical results
do not establish prevalence in operating fleets.
\end{abstract}
\noindent\textbf{Draft status.} This is a focused research draft for critical
assessment of correctness, originality and significance. The public-case
interval below is the presently established interval; no new timetable-scale
certificate is claimed. The draft is not a submission-readiness judgment.

\section{Introduction}
An electric-bus operator assigns mandatory services to vehicles and decides
when to recharge them. Continuous charging offers flexibility, but a vehicle
cannot cover a fraction of a service. Assigning two services to the same bus
can therefore impose an intermediate energy requirement that disappears when
two buses are used. A tariff may reward the resulting change in charging
without compensating the operator for its additional vehicle.

We ask a precise question. A supplier has a convex incremental cost for
serving aggregate fleet demand. A tariff is calculated from a candidate
fleet's load and posted before the operator chooses its complete schedule.
Holding that tariff fixed, does the candidate minimize the operator's bill
plus its intrinsic operating cost? This is an own-load price-support test.
Its participant is one operator controlling the complete fleet. It is not a
strategic model in which that operator anticipates a new tariff after each
alternative it considers.

The connection between price support and convexification is established in
nonconvex market theory. Our question is what a fully specified electric-bus
model can contribute beyond that connection. Full terminal replenishment is
central: an initially full battery must not provide free net energy over a
finite horizon. Our first construction buys exactly the same total energy
under both service assignments, restores every used bus to full, and still
has a strictly positive physical--hull gap. The mechanism is a change in
\emph{when} the fleet purchases energy, coupled with the fixed
cost of a used bus.

The main analysis then asks whether this obstruction is an artifact of an
overly small alternative set. We give an exact connector-capacity and vehicle-
fee characterization, and distinguish two kinds of enlargement. Additional
energy-consuming movements under unchanged equipment and charging access
cannot improve either the physical or hull value in the stated construction.
An added battery type changes the equipment set and can restore support. Its
fee matters: a newly admitted plan cheaper than the old hull optimum need
not restore support because the enlarged hull can improve further. Finally,
an exact replication calculation shows why a small planning gap and a small
whole-owner incentive are different statements.

These are application-specific results in a synthetic model. We do not
claim that nonconvex demand or convex-hull price support is a new concept.
The public timetable evidence is included to state clearly what is known at
a larger scale and what remains unproved. Its bounds do not currently show
that its optimal physical fleet is unsupported. Practical magnitude and
prevalence are separate research questions.

\section{Relation to existing work}
Electric-bus scheduling combines service assignment, battery feasibility and
charging access. Grid-aware multi-depot scheduling, capacitated stations and
scalable route-and-recharge generation are studied by
\citet{wu2022grid,devos2024capacity,parmentier2023fleets}. These operational
problems motivate our use of a complete fleet: charging feasibility is joint
across the selected vehicle duties, especially at a shared connector.
Our question concerns support over alternatives that can change both the
number of used buses and the charging opportunities created by their duties.
Integrated flow and set-partitioning models already combine assignments,
battery feasibility and charging slots \citep{borndorfer2024integrated}.
\citet{wu2022grid} already choose vehicle-to-trip assignments and charging
under grid-related objectives; full overnight charging is among their
assumptions. Endogenous assignments, full charging and grid-aware scheduling
are therefore not, separately or together, sufficient novelty claims here.

\citet{zoltowskaLin2021} formulate aggregation and disaggregation of bus
charging around hourly auction bids using integer models. Their work is a
precedent for discrete bus charging and market participation; our question
additionally allows service assignments and used-bus costs to change under
explicit terminal replacement-energy accounting. More directly,
\citet{wu2019gridinteractive} describe bus service planning coupled to market
clearing and distribution locational marginal prices in a bilevel framework.
Our question is the support of complete physical fleets at a fixed posted
marginal tariff, rather than a first coupling of bus operations to endogenous
electricity prices.

Price-support failure with indivisible operating decisions is established
in electricity-market research. \citet{oneill2005} study prices in markets
with nonconvexities; \citet{gribikHoganPope2007} discuss market-clearing
prices and uplift. Convex-hull pricing provides the relevant convexification
framework \citep{schiro2016,huaBaldick2017}. Linear-price support in
decentralized problems with integralities is also studied through complete
linear descriptions \citep{huembs2022complete}. We use this machinery; we do
not present the support equivalence below as a new general theorem.
Whole feasible participant schedules as columns and their fixed-price
self-scheduling oracles also have established decomposition treatments
\citep{andrianesis2022hull}. The complete-fleet unit ensures the right
physical alternatives in this application; it is not new column-generation
machinery.

Exact EV flexibility aggregation is also established.
\citet{mukhiAbate2023exact} characterize fixed-population charging polytopes
using permutahedra, with prescribed availability and full departure charge.
\citet{mukhiLohoAbate2025exact} give generalized-polymatroid representations
and exact Minkowski aggregation for convex individual flexibility sets.
Our distinction is the joint load--vehicle-cost frontier when service pairing
and the number of used buses change. The load-only projection in our class
is an interval, not a new nonconvex flexibility set. Fixed-matching interval
aggregation and ordered matching are classical ingredients.

Bidirectional charging and transport decisions have also been combined
previously. \citet[Sections 2.1--2.2]{blatiak2022v2g} optimize commercial
fleet trip timing and charging against energy and balancing-service prices.
The public author manuscript of \citet[Section 2]{zhou2025v2g} jointly
chooses bus use, trip connections and partial charging/discharging under
input purchase and export tariffs. Its objective includes vehicle fees,
deadheading, charging, export revenue and degradation, with full-depot
inventory and overnight recharging described in the model and case study.
These features materially overlap our physical setting. Our comparison
uses the author manuscript; imperfectly extracted mathematical glyphs do
not justify detailed claims about its connector or terminal constraints.
The distinction is our exact cost-labelled threshold and own-load support
classification, not a richer operational V2G formulation.
Neither coupling transport to V2G nor using spare buses for grid services
is a novelty claim here.

Signed aggregate flexibility is established as well.
\citet[Section II, Example 1]{altaha2024flexibility} model lossless
bidirectional batteries with initial and final energy requirements and
aggregate a fixed population through Minkowski sums. Our V2G extension
instead classifies the joint load--fleet-cost frontier created by successor
matching and a shared gross-throughput budget. Its thresholds and paid-idle
elimination argument concern that restricted physical class; the
persistence and restoration examples illustrate the classification, rather
than establishing a first signed-aggregation or V2G price-support result.

Lossy storage also has exact coordinate reformulations.
\citet[Theorems 1--3]{altahaBitar2024lossy} establish a piecewise-affine
grid-power/SOC bijection, a convex SOC feasible set for a single store,
and sufficient monotonicity conditions for a convex reformulated objective.
Our common-loss result does not claim these transformations or convexity
conditions as new. It combines service-dependent vehicle counts and
successor choices with shared gross connector use to describe the efficient
joint early-load, terminal-load and vehicle-cost envelope. Its finite
piecewise description is conditional on the stated common-window model,
and excludes dominated points only under explicitly monotone economics.

Aggregation must specify who owns the aggregated decisions. The classical
perspective of \citet{starr1969} and dispatchable-demand results of
\citet{hreinsson2021} explain why aggregate nonconvexities can become small.
Our replication calculation instead isolates a distinction between a
planning-value gap, one owner's absolute fixed-price incentive and the
incentives of independently reserved service pairs. The explicit rates do
not contradict approximate aggregation results.

The comparison concerns what can be decided from an explicit frontier,
rather than whether assignments or storage can be optimized at all:
\begin{center}
\small
\begin{tabular}{p{0.40\textwidth}p{0.52\textwidth}}
\toprule
Established viewpoint & Decision supplied by this classification\\
\midrule
Operational fleet models choose assignments and charging for a specified
objective \citep{wu2022grid,borndorfer2024integrated,zhou2025v2g}.
& Theorem~\ref{thm:sharedprojection} gives every attainable load--vehicle-cost
layer in the stated two-batch class; Corollary~\ref{cor:sharedsupport}
tests physical attainment of its hull optimum.\\
Fixed-population flexibility represents aggregate battery loads
\citep{mukhiAbate2023exact,altaha2024flexibility}.
& Changing successor matches and vehicle counts changes operating-cost
layers even when the load-only set is an interval. The thresholds retain
this distinction.\\
Lossy-storage reformulation relates grid power and inventory
\citep{altahaBitar2024lossy}.
& Theorem~\ref{thm:lossyfrontier} reduces the assignment-dependent efficient
joint load--cost set to at most two segments per cardinality, preserving
planning values under monotone supply costs.\\
\bottomrule
\end{tabular}
\end{center}
These are restricted structural decisions, not claims of a richer operational
model or faster solution of general bus scheduling. Standard matching
integrality, value-function convexity and price-support duality remain the
underlying tools. The capacity, fee and equipment examples explain the
classification's consequences; the separately certified public instance
establishes existence within its declared movement graph, rather than
extending the closed-form theorem to arbitrary timetables.

\section{Complete fleets and the meaning of support}\label{sec:framework}
Fix mandatory services, directed movement opportunities, a charging horizon
and physical equipment. A complete plan $x\in\X$ covers every mandatory
service exactly once, assigns the services to used buses, routes those buses,
and schedules charging subject to battery and shared-resource constraints.
Uncounted buses neither exchange energy nor incur operating cost. When a
variant admits paid idle buses, they count as used participants under that
variant's fleet and access rules. Let $e(x)\in\R^T$ denote net grid energy
by market interval, positive for purchases and negative for exports, and let
$c(x)$ be intrinsic operating cost. Charge-only variants restrict
$e(x)\in\R_+^T$. Each used battery starts full and finishes at the
same required full inventory. Assume the projected image
\[
 S=\{(e(x),c(x)):x\in\X\}
\]
is nonempty and compact. All physical planning, price-response and hull
problems below use exactly this common set.

Let $F:\R^T\to\R\cup\{+\infty\}$ be proper, closed and convex, finite on
all feasible loads. It represents incremental supply cost conditional on
fixed background demand, on the declared feasible signed-load domain when
exports are allowed. Define
\begin{align}
 D&=\min_{x\in\X}\{c(x)+F(e(x))\},\label{eq:D}\\
 \CH&=\min_{(e,c)\in\operatorname{conv}S}\{c+F(e)\},
 \qquad\Delta=D-\CH\ge0.\label{eq:CH}
\end{align}
The hull mixes \emph{complete feasible plans}; it is not a duty-column
relaxation unless equivalence is separately proved. Its mean load need not
be executable as a physical daily schedule. Evaluating supply cost at that
mean also differs from averaging the costs of plans operated on different
days.

For a posted tariff $p\in\R^T$, the complete-fleet response and regret are
\begin{equation}\label{eq:regret}
 V(p)=\min_{y\in\X}\{c(y)+p^\top e(y)\},\qquad
 r(x;p)=c(x)+p^\top e(x)-V(p).
\end{equation}
The posted tariff is fixed across alternatives $y$. A strategic operator
anticipating a tariff reset would face a different objective and needs a
separate settlement model.

\begin{definition}[Weak own-load support]
A physical plan $x$ is supported if there exists an admissible marginal
price $p\in\partial F(e(x))$ for which $r(x;p)=0$. Alternative plans may
tie. If an economic convention restricts the price set, existence is tested
only in that declared subset.
\end{definition}
For differentiable supply we use its gradient. For our quadratic examples
this means the gradient of the unrestricted quadratic formula, including
at a zero load. In charge-only variants, orthant-domain normals are not additional
generation-price choices. With nonsmooth supply, a failed support test at
one chosen multiplier does not exclude support at another admissible price.

For the following classical equivalence, assume either
\[
 \operatorname{ri}\operatorname{conv}e(\X)
 \cap\operatorname{ri}\operatorname{dom}F\ne\varnothing,
\]
or that a finite convex minorant on $\R^T$ agrees with $F$ on all hull
loads. The latter applies to our quadratic. These assumptions ensure an
attained hull-dual price; a restricted economic price set must contain such
a price for the corresponding converse.

\begin{proposition}[Classical support criterion]\label{prop:support}
For every physical $x$ and $p\in\partial F(e(x))$,
\begin{equation}\label{eq:lower}
 r(x;p)\ge c(x)+F(e(x))-\CH\ge\Delta.
\end{equation}
Under the stated qualification, $D=\CH$ if and only if some physical plan
is supported. If $D=\CH$, every physical planner optimum is supported by
every \emph{admissible} attained maximizer of $V(p)-F^*(p)$, where
$F^*(p)=\sup_u\{p^\top u-F(u)\}$.
\end{proposition}
This specializes established nonconvex support and convexification
arguments \citep{oneill2005,gribikHoganPope2007,huaBaldick2017}; the proof
is included in Appendix~\ref{app:support} to make assumptions and quantifiers
explicit. It does not say that every subgradient supports a planner optimum,
or that a price-adjustment algorithm converges. Without a price restriction,
every attained dual maximizer has zero regret and is a subgradient at each
physical planner optimum. An external price restriction adds a separate
admissibility test to that mathematical conclusion.

\begin{remark}[Admissibility is a separate quantifier]
For a single physical point $(e,c)=(0,0)$ and $F(e)=|e|$, both values
are zero and every $p\in[-1,1]$ attains the hull dual. If the declared
admissible price set is $\{0\}$, the plan has admissible support at zero;
the dual maximizer $p=1$ remains mathematically supporting but is inadmissible.
Containing one dual maximizer does not mean containing them all.
\end{remark}

The participant in \eqref{eq:regret} controls all buses. Shared infrastructure
does not automatically define unilateral alternatives for separately owned
vehicles. Section~\ref{sec:ownership} introduces reserved rights before
making an individual-participant comparison.

\section{A fully replenished physical obstruction}\label{sec:construction}
Two mandatory depot-to-depot services consume 15 kWh each: A occupies
$[0,1]$ and B occupies $[2,3]$. At most two identical buses are available,
each with 20 kWh capacity, zero reserve and full initial inventory. Early
charging is available only in $[1,2]$ at 10 kW, both per bus and in total.
Terminal charging is available only in $[3,4]$, using one connector of
power $K\le30$ kW. Every used bus must finish full at hour 4. Charging is
continuous and lossless, with no taper or switching time. Travel cost,
movement consumption and auxiliary load are zero. All prices and costs in
this section are synthetic.

There are exactly two service partitions up to vehicle labels. With one
bus, inventory after A is 5; it needs at least 10 early kWh to complete B.
The early limit therefore fixes load $(10,20)$. With two buses, A's bus
buys $x\in[0,10]$ early; terminal purchases are $15-x$ for A's bus and 15
for B's bus. Each bus returns full. Both partitions buy exactly 30 kWh;
initial inventory cannot subsidize the horizon.

Let $f\ge0$ be the cost of a used bus and define
\begin{equation}\label{eq:supply}
 F(x,z)=4x+\frac{x^2+z^2}{10},\qquad x,z\ge0,
 \qquad g(x)=F(x,30-x)=90-2x+\frac{x^2}{5}.
\end{equation}
The early/late distinction permits different marginal supply costs despite
constant total energy.

\begin{theorem}[Nominal failure]\label{thm:nominal}
For $K=30$ and $f=7$,
\[
 D=97,\qquad \CH=\frac{7591}{80}=94.8875,\qquad
 \Delta=\frac{169}{80}>0.
\]
The unique physical planner load is $(10,20)$ with one bus. Its own-price
regret at $p=(6,4)$ is 13.
\end{theorem}

The one-bus objective is $7+g(10)=97$; the two-bus branch minimizes
$14+g(x)$ at $x=5$, giving 99. The hull instead mixes one bus at early
energy 10 with two buses at early energy 0. Its cheapest operating cost at
mean early energy $x$ is $14-0.7x$. Thus its objective is
\[
 14-0.7x+g(x)=\frac{7591}{80}+\frac15(x-27/4)^2.
\]
The mixture weights are $27/40$ and $13/40$. Both components are physical;
the mean operating cost $371/40$ is a convexification. Running those plans
on different days would have average physical cost 99.275, not 94.8875.

At the physical optimum's marginal price $(6,4)$, its bill plus operating
cost is 147. The two-bus zero-early alternative costs 134 at that fixed
tariff. This 13-unit incentive changes the entire fleet. By contrast, the
hull price $(107/20,93/20)$ makes the two hull components tie, but is not
the marginal supply price of the one-bus physical load. The distinction
explains why a hull price and an own-load tariff answer different questions.

\subsection{An exact connector and fee characterization}
Define $\clamp(v;\ell,u)=\min\{u,\max\{\ell,v\}\}$. The next result
describes the entire declared parameter domain, including its boundaries.

\begin{theorem}[Connector capacity and vehicle fee]\label{thm:capacity}
For $0\le K\le30$ and $f\ge0$, the model is infeasible if $K<20$.
If $K=20$, $D=\CH=f+90$; both service partitions tie only at $f=0$.
For $K>20$, put
\[
 \ell=30-K,\quad d=K-20,\quad t=\max\{5,\ell\},\quad
 h=\clamp\left(5+\frac{5f}{2d};\ell,10\right).
\]
Then
\begin{align}
 D&=\min\{f+90,\ 2f+g(t)\},\label{eq:capacityD}\\
 \CH&=2f-\frac{f(h-\ell)}d+g(h).\label{eq:capacityCH}
\end{align}
For $f>0$, the gap is strictly positive exactly when
\begin{equation}\label{eq:capacityregion}
 \max\left\{0,\frac25d(\ell-5)\right\}<f<2d.
\end{equation}
For $f=0$ the gap is zero. Both physical planner partitions are retained
when $f=90-g(t)$.
\end{theorem}

The proof is a physical projection and one-dimensional convex minimization,
not a finite parameter scan (Appendix~\ref{app:capacity}). In particular,
with $f=7$ the gap is zero for $20\le K\le47/2$ and positive for
$47/2<K\le30$. Extra charging capacity can create a price-support failure
by enlarging the alternative set; it does not increase the planner's minimum
cost or show that installing a connector worsens welfare.

The full price-response value is
\begin{equation}\label{eq:fullresponse}
 V(p)=\min\{f+10p_1+20p_2,\;
 2f+\ell p_1+(30-\ell)p_2,\;2f+10p_1+20p_2\}.
\end{equation}
An affine objective attains its minimum on each physical segment at an
endpoint. If $p_1=p_2$, its whole two-bus interval is minimizing within
that partition and must be retained whenever its bill also ties globally.
This includes the zero-fee boundaries; an endpoint representative alone
would omit part of the response set.

\subsection{Positive reserve and charging losses}
The nominal early limit is binding: adding reserve or losses without
changing power can remove the one-bus branch. This is a feasibility issue,
not disappearance of the pricing mechanism. A nearby construction uses
reserve 1 kWh, efficiency $19/20$, early power 12 kW and one terminal
connector of 30 kW, keeping the service energies, capacity 20 and fee 7.
The one-bus state sequence is
\[
 20\longrightarrow5\longrightarrow16\longrightarrow1
 \longrightarrow20.
\]
It buys $220/19$ grid kWh early and 20 terminally. A physical two-bus
endpoint buys $30/19$ early and 30 terminally; terminal sessions of
$9/19$ and $10/19$ hour fit serially.

\begin{proposition}[A nondegenerate neighborhood]\label{prop:reserve}
For these modified parameters,
\[
 D=\frac{38527}{361},\qquad
 \CH=\frac{2987911}{28880},\qquad
 \Delta=\frac{94249}{28880}>0.
\]
With terminal power 30 fixed, a positive gap persists in an open
neighborhood of the stated reserve, efficiency and early power.
\end{proposition}
Appendix~\ref{app:reserve} states the strict inequalities defining this
neighborhood. The result remains conditional on constant efficiency,
zero switching time and the absence of charging taper.

\section{Does failure survive additional alternatives?}\label{sec:expansion}
Let $X$ be an old set of complete plans, $Y$ a set of additions, and
$Z=X\cup Y$. Use the same supply function and compact projected-image
assumptions as before, with attained hull-dual prices where needed. Write
$D_A,\CH_A,V_A$ for the values associated with a set $A$. Suppose
$\Delta_X>0$. The following distinction is essential: a retained old plan
cannot be rescued by more deviations, but a newly admitted plan can become
a supported optimum.

\begin{proposition}[Expansion and an old dual certificate]\label{prop:expansion}
For every old $x$ and unchanged own-load price $p$,
\[
 r_Z(x;p)\ge r_X(x;p)>0.
\]
Erasure of global failure therefore requires a new supported plan and,
necessarily, $D_Y\le\CH_X$. A sufficient condition for survival is
$\min\{D_X,D_Y\}>\CH_X$.

Fix an attained old hull-dual optimizer $p^*$. If every added plan satisfies
\begin{equation}\label{eq:dualsafe}
 c(y)+(p^*)^\top e(y)\ge V_X(p^*),
\end{equation}
then $\CH_Z=\CH_X$. Under this condition, failure is erased exactly when
some added plan attains equality in \eqref{eq:dualsafe} and also satisfies
$p^*\in\partial F(e(y))$. If \eqref{eq:dualsafe} holds with a uniform
margin $\varepsilon>0$, then
\[
 \Delta_Z\ge\min\{\Delta_X,\varepsilon\}>0.
\]
\end{proposition}

These are elementary consequences of the classical dual certificate
(Appendix~\ref{app:expansion}). They provide a useful way to ask whether
a restricted physical mechanism survives genuinely richer alternatives.
With strictly convex supply, supply equality at $p^*$ fixes the old
hull-optimal load; fleet equality then fixes its operating cost. Under
\eqref{eq:dualsafe}, erasure means physically implementing that projected
hull optimum. Without \eqref{eq:dualsafe}, the hull can move, and comparison
with its old value is not enough.

\subsection{Movement additions under unchanged resources}
In the nominal model of Section~\ref{sec:construction}, retain all original
plans and allow additional movements consuming nonnegative energy and
incurring nonnegative extra operating cost. Preserve the two services,
at most two buses with 20 kWh batteries, fee 7, full initial and terminal
inventories, the early 10 and terminal 30 kW power limits, charging windows and access.
There is no energy recovery, inventory transfer or new charging opportunity.

\begin{theorem}[Movement-only stability]\label{thm:movement}
Every feasible added plan is componentwise dominated in purchased-energy
and operating-cost space by an old physical plan. Consequently
\[
 D_Z=97,\qquad\CH_Z=\frac{7591}{80},\qquad
 \Delta_Z=\frac{169}{80}.
\]
\end{theorem}

To see the physical reason, any added plan buys $30+w$ total energy,
$w\ge0$, and at most 10 early kWh. A two-bus plan with early purchase $u$
is dominated by the old two-bus plan $(u,30-u)$ of cost 14. A one-bus
plan must still obtain at least 10 early kWh to complete both services;
the early limit forces $u=10$, and the old one-bus plan $(10,20)$ of
cost 7 dominates it. The same replacement works inside every hull mixture
because $F$ is coordinatewise nondecreasing. The theorem therefore preserves
both values, not merely positivity of their difference.

This is an unchanged-resource result. The original no-movement construction
already includes both possible service partitions and every continuous
charging choice. A new projected point cannot be called an overlooked
feasible duty without explaining which physical assumption admits it.

\subsection{An explicitly enlarged battery choice}
Now change the equipment. Keep buses with 20 kWh batteries and fee 7, and add buses
of capacity $B\in[20,30]$ and fee $q\ge7$. Two of each type are available,
but at most two buses total are used. Every used bus starts and ends at
its own full capacity. All service and charging constraints are unchanged.

Put $a=30-B$. The new one-bus branch buys $x\in[a,10]$ early and
$30-x$ terminally, with operating cost $q$. Every two-bus type assignment
permits the old interval $[0,10]$ and costs at least 14; the old--old
assignment therefore dominates mixed and new--new assignments in projected
economic comparisons. Their physical records and equal-fee ties remain
part of the complete alternative set.

\begin{theorem}[Battery and fee conditions for erasure]\label{thm:battery}
For $20\le B\le30$ and $q\ge7$, let $a=30-B$ and
\[
 A(a)=14+2a-\frac25a^2,\qquad
 C(a)=7+(10-a)\left(\frac25a-2\right).
\]
The enlarged model has a supported physical plan exactly when
\begin{equation}\label{eq:batteryregion}
 \bigl(a\le5\ \hbox{and}\ q=7\bigr)
 \quad\hbox{or}\quad
 \bigl(a>5\ \hbox{and}\ 7\le q\le\min\{A(a),C(a)\}\bigr).
\end{equation}
The inequalities include weak support at ties. At fee $q=7$, this becomes
\begin{equation}\label{eq:batterythreshold}
 B\ge\frac{55-\sqrt{95}}2\simeq22.6266.
\end{equation}
Outside \eqref{eq:batteryregion}, the physical--hull gap is positive.
\end{theorem}

The new branch's physical minimizer is $t=\max\{5,a\}$. If $a\le5$,
its marginal prices at $t=5$ are equal; the new plan minimizes its bill
against the cheaper old one-bus plan only at fee 7. If $a>5$, the two
inequalities in \eqref{eq:batteryregion} compare it respectively with
the old zero-early two-bus plan and old one-bus plan. The proof in
Appendix~\ref{app:battery} checks the entire continuous opportunity set.

The case $B=25$ shows why a new plan's physical cost is not the whole
story. For $q\ge7$,
\begin{align}
 D_Z&=\min\{97,85+q\},\label{eq:B25D}\\
 \CH_Z&=
 \begin{cases}
 85+q-(q-7)^2/20,&7\le q\le21/2,\\
 7591/80,&q\ge21/2.
 \end{cases}\label{eq:B25CH}
\end{align}
At fee 7, the new one-bus load $(5,25)$ has total cost 92 and is supported
at $(5,5)$. At fee 8, its physical cost 93 is below the \emph{old} hull
value 94.8875, but the \emph{new} hull value is 92.95. Its regret is 1 and
the gap remains $1/20$. Thus the necessary condition $D_Y\le\CH_X$
in Proposition~\ref{prop:expansion} is not sufficient.

At fee 12, the old one-bus load $(10,20)$ and new one-bus load $(5,25)$
both have physical cost 97. Their own-price regrets are 13 and 5,
respectively; $\CH_Z=7591/80$. Reporting a single selected planner
optimum would hide this incentive difference. The new equipment can erase
failure, but its presence alone does not do so.

\section{Aggregation and the owner of the deviation}\label{sec:ownership}
The homogeneous specialization of Section~\ref{sec:heterogeneous} provides
a separate way to distinguish a planning-value gap from a participant's
incentive. Replicate the nominal services to $n$
simultaneous copies of A and $n$ copies of B. Retain the per-bus physical
parameters and fee 7. Provide $n$ early connectors of power 10 kW and $n$
terminal connectors of power 30 kW. Supply scales with the market:
\begin{equation}\label{eq:scaledF}
 F_n(E,L)=4E+\frac{E^2+L^2}{10n}.
\end{equation}
This scaling keeps the marginal-price range bounded; it is an assumption
of the family, not an asymptotic claim for an unchanged supply curve.

If $m$ buses each cover an A--B pair, the remaining services use $n-m$
A-only and $n-m$ B-only buses. Operating cost is $14n-7m$, early energy
is in $[10m,10n]$, and terminal energy is $30n-E$. All terminal charges
fit: each paired bus takes one connector, while each unpaired A/B pair
shares one connector serially. Thus the projected opportunities used below
have explicit physical realizations.
All A-to-B crossmatchings are admitted in this aggregate model. Identical
energies make their projected outcomes depend only on $m$; fixed pair duties
are not being imposed here. In the threshold notation, $r_m=10m$.

\begin{theorem}[Planning gap and whole-owner regret]\label{thm:replication}
The hull value is $\CH_n=7591n/80$. Every physical planner optimum takes
an integer $m$ nearest to $27n/40$, retaining both choices at a tie.
Writing $\delta=m-27n/40$,
\begin{equation}\label{eq:scaledgap}
 \Delta_n=\frac{20\delta^2}{n}\le\frac5n.
\end{equation}
For one owner controlling all services and buses, the own-price regret is
\begin{equation}\label{eq:scaledregret}
 r_n=
 \begin{cases}
 m\,40\delta/n,&\delta\ge0,\\
 (n-m)(-40\delta/n),&\delta<0.
 \end{cases}
\end{equation}
In particular, along $n=40k+1$,
\[
 \Delta_n=\frac{169}{80n}\longrightarrow0,\qquad
 r_n=\frac{351}{40}+\frac{169}{40n}
 \longrightarrow\frac{351}{40}.
\]
For multiples of 40, both quantities are zero.
\end{theorem}

Approximate convexity of an aggregate does not by itself bound a composite
owner's absolute incentive to replace its entire aggregate. This is consistent
with aggregation perspectives \citep{starr1969,hreinsson2021}. In this
family regret per bus and regret relative to physical cost vanish; the
nonvanishing statement concerns the whole owner's \emph{absolute} fixed-price
gain along a declared subsequence.

Now divide ownership among $n$ operators, each responsible for one A/B
pair, with a reserved early connector and a reserved terminal connector.
Each may use one or two buses and choose its two-bus early charge in
$[0,10]$. Those reserved rights ensure unilateral deviations remain feasible
with every other participant's plan held fixed.
An aggregate optimum may pair A and B copies across these ownership labels.
Before making an individual comparison, implement its same projected outcome
by assigning $m$ owners a one-bus plan and the remaining $n-m$ owners a
two-bus plan on their own reserved service pairs. Identical service copies
and the explicit connector assignment make this possible with unchanged
aggregate load, operating cost and marginal price. This is a reimplementation
of the economic outcome, not a claim that a cross-matched literal schedule
already respects the newly imposed ownership rights.
This step is specific to identical services. In the heterogeneous example
of Section~\ref{sec:heterogeneous}, fixed pairs with energies $(5,14)$
and $(15,15)$ require 10 early kWh when both use one bus, whereas the
crossed complete assignment requires only 9. The latter outcome cannot
be transferred to that fixed-pair institution by relabeling. No individual
regret bound for heterogeneous fixed ownership is asserted below.

\begin{corollary}[Reserved participants]\label{cor:individual}
For any physical planner optimum reimplemented in the reserved-pair
institution as above, each reserved pair's regret at its marginal price
is at most $20/n$. The regrets sum to
\eqref{eq:scaledregret}. Along $n=40k+1$, each dissatisfied one-bus owner
can gain only $13/n$.
\end{corollary}
The proof is the same per-pair bill comparison
(Appendix~\ref{app:replication}). At $n=20$ the two tied physical planner
optima with $m=13$ and $m=14$ have whole-owner regrets 7 and 14,
respectively. Neither the whole-owner calculation
nor the reserved-pair calculation specifies an institution for unreserved
shared chargers or strategic anticipation of a tariff reset.

\subsection{Which uniform controls make a small gap informative?}
The preceding example does not imply that planning error and incentives are
unrelated. The selected price matters in the classical link between lost
opportunity cost and duality gaps \citep{huaBaldick2017}. The following
elementary convex-analysis bound makes the relevant uniform controls explicit;
it is an explanatory consequence, not a new general pricing principle.

\begin{proposition}[Gap and own-price regret under uniform curvature]
\label{prop:regretbound}
Let the projected complete-plan set $S=\{(u(x),c(x)):x\in\X\}$ be nonempty
and compact, and let $K$ be its load projection, of diameter $R_K$.
Suppose $F$ is differentiable, $\mu$-strongly convex and has
$\Lambda$-Lipschitz gradient on a relative neighborhood of
$\operatorname{conv}K$, where $0<\mu\le\Lambda$.
Use one fixed Euclidean norm on its affine hull.
For any physical planner optimum $x$, at $p=\nabla F(u(x))$,
\[
 \Delta\le r(x;p)\le
 \Delta+\Lambda R_K\sqrt{2\Delta/\mu}.
\]
If loads lie in a proper affine subspace, gradients and norms can be taken
tangentially: normal price components add the same amount to every bill.
Economic price admissibility is a separate qualification.
\end{proposition}
\begin{proof}
Let $(u^*,c^*)$ minimize $c+F(u)$ on $\operatorname{conv}S$ and put
$q=\nabla F(u^*)$. Hull optimality and linear minimization give
$V(q)=c^*+q\cdot u^*$. With the Bregman difference
$B_F(u,u^*)=F(u)-F(u^*)-q\cdot(u-u^*)$, this yields the exact identity
\[
 c(x)+q\cdot u(x)-V(q)=\Delta-B_F(u(x),u^*)\ge0.
\]
Consequently $\|u(x)-u^*\|\le\sqrt{2\Delta/\mu}$.
For this fixed plan, regret is the maximum of affine price functions with
slopes $u(x)-u(y)$, so it is $R_K$-Lipschitz in price. Transfer from $q$ to
$p$, use gradient smoothness and discard the nonpositive Bregman term to
obtain the upper bound. The lower bound is the tangent inequality from
Proposition~\ref{prop:support}.
\end{proof}

Uniformly bounded $R_K,\Lambda$ and a uniform positive lower bound on
$\mu$ therefore make planner regret vanish as $\Delta\to0$. In our
replication family, however,
$\mu_n=\Lambda_n=1/(5n)$ and $R_{K,n}=10\sqrt2\,n$ in two-load
coordinates. The displayed upper becomes $\Delta_n+40|\delta|$.
Indeed the fixed-$q$ score gap is zero there, so retaining its Bregman
term sharpens this to $40|\delta|$. The exact nonvanishing whole-owner
regret is thus consistent with the bound. Arbitrary nonoptimal plans require
their own excess $c(x)+F(u(x))-\CH$ in place of $\Delta$; nonsmooth supply
or vanishing curvature needs a different argument.

\section{What the public timetable does and does not establish}\label{sec:public}
The analytical construction proves possibility, not operational frequency.
The public case uses the Hildenbrand timetable from the Sistig data
\citep{sistig2025dataset,sistig2025electrification}: 37 mandatory services
and the complete directed deadhead matrix. Depot 15 and depot 16 define
separate one-depot modifications of the original network. Source timetable
offsets, including next-day services, are retained; source data do not
provide the synthetic energy and cost assumptions below.

Service consumption is 1.20 kWh/km plus 16 kWh per service hour; directed
deadheads use 0.96 kWh/km plus 16 kWh per travel hour, and explicit off-depot
waiting uses 16 kWh/h. Depot-parked auxiliary draw is zero. A used bus has
400 kWh usable inventory, zero reserve, unit efficiency and 360 kW acceptance.
One 360 kW connector is shared. Every used bus starts full and returns full
by a synthetic 30:00 service-block boundary. There is no taper or switching
time. Monetary movement cost is zero; movement time and energy remain
physical constraints. The used-bus fee is 100. This finite block is not a
certificate of repeating daily operation.

The initial flat tariff is 0.2 per kWh. Because the complete-plan objective
is linear, its physical and hull optima have equal value regardless of the
discrete assignments. A native numerical flat optimum near 408.53 at depot 15
is tolerance-qualified and does not identify the ideal stored-input rational
optimum. Exact rational evidence instead provides a two-bus physical witness
and a cardinality-gated path-cover relaxation, giving outward bounds
\[
 D^{\rm flat}_{15}\in[404.92,408.54].
\]
An exact one-bus obstruction establishes minimum fleet size two. Depot 16
has an exact lower bound 414.39, but its reported two-bus upper witness is
numerical; the resulting interval must be labeled mixed evidence.

On the same finite depot 15 ideal model, let
\begin{equation}\label{eq:publicF}
 F(e)=\sum_{t=1}^{30}\left(0.2e_t+\frac{e_t^2}{1800}\right).
\end{equation}
The displayed coefficients denote the original stored binary values in the
exact certificates.\footnote{Precisely, the linear coefficient is
$3602879701896397/18014398509481984$ and the quadratic Hessian coefficient is
$5124095576030431/4611686018427387904$; the quadratic term is half the latter
times $e_t^2$. The familiar decimal expression displays their intended scale.}
A new exact two-bus witness and an independently reconstructed rational
lower certificate for a per-leg perspective energy-flow relaxation
give the presently established outward bounds
\begin{equation}\label{eq:publicbounds}
 434.17\le\CH_{15}\le D_{15}\le441.54,\qquad
 0\le D_{15}-\CH_{15}\le7.37.
\end{equation}
The new exact lower is approximately $434.17948$, rounded downward here;
its validity follows from a supply tangent, sign-valid rational row
multipliers and finite-bound compensation for stationarity residuals.
It improves the earlier energy-floor lower of 429.22. The continuous
energy-flow relaxation is not asserted to equal the complete-fleet hull.
A separate native integer solve returned OPTIMAL with objective
$441.532912\ldots$, but its candidate failed strict physical replay and
did not itself improve the exact upper or establish ideal-model optimality.
The native status and the independently checked rational lower are
different evidence. The new admitted witness has cost
$442.123049\ldots$. A time-limited classical MIQP supplied its routes and
charging proposal; two terminal energy changes smaller than $7\cdot10^{-12}$
kWh restored exact full inventories. A separate rational event replay then
verified all services, directed movements, battery bounds, shared connector
use and full replenishment. The original unadmitted numerical proposal is
retained separately. No route or nonterminal charging allocation changed.

For any fleet with at least three buses, its operating cost and only the
linear mandatory-service-energy component already exceed $477.82$. Movement
consumption and the quadratic cost are nonnegative. Such a fleet cannot
improve the new physical upper. Together with the one-bus obstruction, this
establishes exactly two buses at any physical planner optimum. It does not
exclude larger-fleet components from arbitrary hull mixtures or from a
price-response problem at another tariff.

Separately, a second strictly admitted physical plan and the protected witness
give a feasible complete-plan mixture with value $441.600681\ldots$, hence
$\CH_{15}\le441.61$. Its negative initial directional derivative proves that
the protected witness is unsupported at its own marginal prices. A mixture
is not a dispatch, and an unsupported incumbent is not proof that an optimal
physical fleet is unsupported. A subsequent fixed-route charging QP returned
the same lower numerical value $441.532912\ldots$, but its first exact dual
recovery failed on a multiplier sign; the conditional global own-price
checks were therefore unattempted. A distinct rational phase-I method then
constructed a new exact primal, admitted it by independent physical replay,
and found sign-valid multipliers certifying its fixed-route charging optimum.
This supplies the stronger physical upper in \eqref{eq:publicbounds}; it also
supersedes the earlier mixture upper as a bound on $\CH$. The original failure
remains recorded. Fixed-route optimality still does not exclude a profitable
complete-fleet deviation or certify global marginal-price support.
The bounds neither determine
$D_{15}$ or $\CH_{15}$ nor establish a positive gap. No new public sign
certificate is asserted in this draft.

There are distinct valid ways to close this question. An exact physical
witness $x$ and a globally valid linear-response lower bound matching
$c(x)+\nabla F(e(x))^\top e(x)$ prove support and hence $D=\CH$.
Alternatively, a physical lower bound strictly above a feasible mixture's
hull upper bound proves positive gap. A globally certified hull optimizer
can also prove failure through nonimplementability: strict convexity in
\eqref{eq:publicF} fixes its load and operating cost, and a hull operating
cost not in $100\mathbb Z_+$ cannot be the cost of a physical fleet.
All these routes require a global certificate or a complete infeasibility
argument. A finite plan bank, incumbent regret, or a small native solver
gap does not settle the sign of the ideal-model gap.


\section{Implications and limits}
The examples identify a conditional failure of a specified posted-tariff
institution. The supply function is convex; the indivisibility lies on the
fleet side. Continuous energy decisions and terminal energy accounting do
not eliminate that indivisibility. Conversely, an indivisible fleet need
not lack supporting marginal prices: the exact parameter regions include
zero-gap cases, and an added battery choice can physically implement a
supported load. The correct claim is a possibility with characterized
conditions, not inevitable failure of marginal tariffs.

The expansion results are deliberately explicit about resource changes.
The movement-only stability theorem preserves charging access, equipment,
service requirements and nonnegative movement consumption. It says nothing
about a new charging location, regenerative energy, altered travel energy,
service cancellation or a more capable battery. Those changes can alter the
relevant energy projection. The equipment result studies one such change
directly rather than describing an arbitrary projected point as an omitted
feasible duty of the original model.

The participant definition is equally consequential. A whole-fleet deviation
can replace every service assignment. Independent operators need compatible
resource rights before their unilateral alternatives are well defined.
The replication theorem uses reserved service pairs and connectors to make
that comparison exact. It does not infer incentives for independently owned
buses competing for the same unreserved depot connector.

All analytical models omit stochastic delays, battery degradation,
temperature effects, charging taper and switching time. Monetary and supply
coefficients are synthetic. The positive-gap neighborhood admits a reserve,
constant losses and one connector, but remains conditional on those
abstractions. The analytical formulas and exact implementation comparisons
support the declared mechanisms, not estimates of an operator's savings or
an assessment of a deployed market design.

The strongest present contribution is therefore a focused analytical study:
a fully replenished physical realization, exact stability and erasure
regions, and a participant-scale distinction. Whether these results are
sufficiently original and significant relative to the closest prior work
requires critical assessment. A resolved public support certificate and
independent timetable evidence would strengthen the case; neither should be
implied by the current bounds. No learning advantage is established here.

\section{Conclusion}
An electric-bus fleet can lack marginal-price support even when charging is
continuous and every used battery is fully replenished. The two-service
model permits an exact answer as connector capacity, vehicle fees and
battery options change. Energy-consuming movement additions preserve the
obstruction under unchanged resources; a sufficiently capable battery can
remove it at an appropriate fee. Neither a cheaper newly admitted plan nor
a vanishing planning gap alone resolves the relevant incentive question.
The public case remains a certification problem rather than empirical
confirmation of the mechanism.

\section*{Evidence and reproducibility}
The analytical claims concern the stated ideal models. Exact rational
implementation comparisons use independently formulated physical enumeration
and analytical formulas; agreement is an internal cross-check, not external
replication. The connector/fee qualification contains twenty fixed cases,
and the support-extension qualification fourteen fixed cases, with explicit
boundary and corruption controls. These cases qualify implementations;
the proofs establish the continuous-domain statements. Archived outcome and
resource records distinguish exact arithmetic from tolerance-qualified
native computations. No new numerical execution was performed to author
this draft, and no public optimal-fleet support conclusion is drawn from
the finite plan bank or solver statuses.

\appendix
\section{Support, physical projection and parameter proofs}
\subsection{The classical support criterion}\label{app:support}
\begin{proof}[Proof of Proposition~\ref{prop:support}]
For any hull point $(e,c)$, linear minimization gives
$c+p^\top e\ge V(p)$. If $p\in\partial F(e(x))$, then
$F(e)\ge F(e(x))+p^\top(e-e(x))$. Adding and minimizing over the hull
gives
\[
 \CH\ge V(p)+F(e(x))-p^\top e(x)
       =c(x)+F(e(x))-r(x;p).
\]
This proves \eqref{eq:lower}. If $r(x;p)=0$, equality of physical and hull
values follows immediately.

Conversely let $D=\CH$. Compactness and lower semicontinuity give a
physical minimizer. The relative-interior qualification gives strong
Fenchel duality and an attained maximum
$\CH=V(p^*)-F^*(p^*)$. Under the alternative finite-minorant condition,
apply the same duality theorem to the finite minorant. It agrees with $F$
on hull loads and its conjugate is at least $F^*$. The corresponding
price gives a dual bound for $F$ at least $\CH$; weak duality forces
equality, hence attainment for $F$ as well.

For any physical minimizer $x$ and any attained hull-dual maximizer $p^*$,
\[
 \underbrace{c(x)+(p^*)^\top e(x)-V(p^*)}_{\ge0}
 +\underbrace{F(e(x))-(p^*)^\top e(x)+F^*(p^*)}_{\ge0}
 =D-\CH=0.
\]
Both terms vanish. Fenchel equality gives $p^*\in\partial F(e(x))$;
the first term gives support. For restricted economic prices the converse
requires an attained dual price in that subset, as stated.
Every dual maximizer satisfies the displayed unrestricted algebraic
conclusion; only the admissible maximizers give support under a restricted
economic definition.
\end{proof}

\subsection{Physical projection and the connector/fee region}\label{app:capacity}
\begin{proof}[Proofs of Theorems~\ref{thm:nominal} and~\ref{thm:capacity}]
With one bus, the state after A is 5 and the state after B is $x-10$.
Early charging has $x\le10$, so feasibility forces $x=10$. Terminal
replenishment is 20 and requires $K\ge20$. With two buses, A's bus has
terminal demand $15-x$, B's bus has demand 15, and $0\le x\le10$.
Serial charging at power $K$ fits the terminal hour exactly when
$30-x\le K$. These conditions also satisfy each battery's individual
headroom and inventory bounds. Thus for $K\ge20$ the two-bus early
interval is $[\ell,10]$, $\ell=30-K$; for $K<20$ neither partition
is feasible. At $K=20$ both projected loads are $(10,20)$ and their
costs differ only by $f$.

For $K>20$, strict convexity of $g$ gives the two-bus minimum at
$t=\max\{5,\ell\}$. The one-bus minimum is the fixed point 10,
proving \eqref{eq:capacityD}. Projecting complete plans to early energy
and operating cost gives the segment $\{(x,2f):\ell\le x\le10\}$
and the point $(10,f)$. Their convex hull has lowest cost
\[
 c_{\min}(x)=2f-\frac{f(x-\ell)}{10-\ell},\qquad\ell\le x\le10.
\]
The derivative of $c_{\min}(x)+g(x)$ vanishes at
$5+5f/(2d)$, $d=10-\ell$, proving \eqref{eq:capacityCH}.

For $f>0$, an interior hull minimizer has lower operating cost than every
two-bus physical plan at the same early load. Strict convexity excludes
the distinct one-bus point from attaining its objective. Therefore the
gap is positive exactly when $\ell<h<10$. Equivalently
$\tfrac25d(\ell-5)<f<2d$, together with $f>0$, which gives
\eqref{eq:capacityregion}. At an endpoint the physical endpoint attains
the hull minimum. For $f=0$, $h=t$, and the two-bus physical plan attains
the hull. Equality of the two physical branch minima gives exactly
$f=90-g(t)$; within each branch strict convexity fixes its minimizing
load. The affine-response formula \eqref{eq:fullresponse} follows by
minimizing over every point of those physical segments.

For $K=30,f=7$, the lower hull objective is
$104-27x/10+x^2/5=7591/80+(x-27/4)^2/5$. Physical branch values are 97
and 99. The one-bus gradient is $(6,4)$, and the complete two-bus
response at that price is the zero-early endpoint. Bills 147 and 134 give
regret 13. The mixture and operating-cost values in the text follow from
weight $27/40$ on the one-bus point.
\end{proof}

\subsection{A reserve/loss neighborhood}\label{app:reserve}
\begin{proof}[Proof of Proposition~\ref{prop:reserve}]
For reserve $0\le r\le5$, constant efficiency $0<\eta\le1$, early
power $P$, and terminal power $K$, define grid-energy quantities
\[
 T=30/\eta,\quad \ell=\max\{0,T-K\},\quad
 u=\min\{P,15/\eta\},\quad
 h=\max\{\ell,(10+r)/\eta\}.
\]
The two-bus early interval is $[\ell,u]$; the one-bus interval is
$[h,u]$ if nonempty. The lower terminal-capacity condition is
$T-x\le K$. For the one-bus plan, its state after B is $\eta x-10$,
so reserve imposes $x\ge(10+r)/\eta$. Charging only A's bus early
realizes every interval point; terminal demands fit serially, and the
individual bounds follow from the inventory headroom 15 and the stated
reserve range.

Write $G_T(x)=F(x,T-x)=G_T(m)+k(x-m)^2$, where $k=1/5$ and
$m=T/2-10$. Suppose
\begin{equation}\label{eq:strictreserve}
 \ell<m<h<u,\qquad f>0,\qquad f<2k(h-\ell)(h-m).
\end{equation}
On $[\ell,h]$ the hull's lowest operating cost is
$2f-f(x-\ell)/(h-\ell)$, obtained from the two-bus endpoint $\ell$
and one-bus endpoint $h$; on $[h,u]$ it is $f$. Put $d=h-\ell$.
The hull minimizer $x^*=m+f/(2kd)$ lies in $(m,h)$. Physical and hull
values are
\begin{align*}
 D&=G_T(m)+\min\{f+k(h-m)^2,2f\},\\
 \CH&=G_T(m)+2f-\frac{f(m-\ell)}d-\frac{f^2}{4kd^2}.
\end{align*}
Subtracting gives
\[
 \Delta=\min\left\{k(h-x^*)^2,
 \frac{f(m-\ell)}d+\frac{f^2}{4kd^2}\right\}>0.
\]
At $(r,\eta,P,K,f)=(1,19/20,12,30,7)$,
\[
 (\ell,m,h,u,d)=(30/19,110/19,220/19,12,10).
\]
In particular $u-h=8/19>0$ and $2kd(h-m)=440/19>7$.
Rational substitution gives the stated values. All inequalities in
\eqref{eq:strictreserve} are strict and the parameter functions are
continuous locally. They therefore persist in an open neighborhood of
$(r,\eta,P)$ with $K=30$ fixed.
\end{proof}

\section{Expansion and equipment proofs}
\subsection{The old-plan and dual-certificate statements}\label{app:expansion}
\begin{proof}[Proof of Proposition~\ref{prop:expansion}]
Set inclusion gives $V_Z(p)=\min\{V_X(p),V_Y(p)\}\le V_X(p)$, so
old-plan regret cannot decrease. By Proposition~\ref{prop:support}, no old
plan can become supported. If failure is erased, $D_Z=\CH_Z\le\CH_X<D_X$;
thus a new plan attains physical optimum and $D_Y\le\CH_X$. If instead
$\min\{D_X,D_Y\}>\CH_X$, then $D_Z>\CH_X\ge\CH_Z$, proving survival.

Under \eqref{eq:dualsafe}, $V_Z(p^*)=V_X(p^*)$, so
\[
 \CH_Z\ge V_Z(p^*)-F^*(p^*)=\CH_X\ge\CH_Z.
\]
For any new physical plan,
\begin{align*}
 c(y)+F(e(y))-\CH_X
 ={}&[c(y)+(p^*)^\top e(y)-V_X(p^*)]\\
 &+[F(e(y))-(p^*)^\top e(y)+F^*(p^*)].
\end{align*}
Both brackets are nonnegative. Equality holds exactly under the fleet and
supply equalities stated in the proposition. Compactness ensures that if
$D_Z=\CH_X$, some physical plan attains that equality; it cannot be old.
With a uniform fleet margin $\varepsilon$, every new physical plan lies
at least $\varepsilon$ above $\CH_X$, while every old one lies at least
$\Delta_X$ above it. This gives the quantitative bound. Strict convexity
of $F$ makes the minimizer of $F(e)-(p^*)^\top e$ unique on its convex
domain whenever attained, proving the projected-point interpretation.
\end{proof}

\begin{proof}[Proof of Theorem~\ref{thm:movement}]
Each added plan buys total energy $30+w$, $w\ge0$, because its initial
and final battery inventories agree. Its early purchase $u$ is at most 10.
For two used buses, the old $(u,30-u)$ plan of cost 14 is feasible and
componentwise no larger in load or cost. For one bus, completion of both
services by hour 3 requires $20+u\ge30$ even before any additional
nonnegative movement consumption. Thus $u=10$, and the old one-bus
plan dominates it. No added feasible plan can reduce the physical
objective because $F$ is coordinatewise nondecreasing.

Every hull point is a finite convex combination of physical points.
Replace each added component by its old dominator; its mixture has no
larger cost or load in any coordinate. Consequently it has no larger
convex supply objective. This gives $\CH_X\le\CH_Z$, and inclusion
gives the reverse inequality. The same inclusion/dominance argument
gives $D_X=D_Z$. This proof does not require the movement timing to be
favorable: it concerns every added plan that remains feasible.
\end{proof}

\subsection{Battery enlargement and fee thresholds}\label{app:battery}
\begin{proof}[Proof of Theorem~\ref{thm:battery} and
\eqref{eq:B25D}--\eqref{eq:B25CH}]
The one-new-bus state sequence imposes $x\ge30-B=a$ and $x\le10$.
Every point of $[a,10]$ is feasible and fully replenished. With two
buses, the A bus's headroom after its service is 15 for either type;
B's bus starts full before service and cannot buy early energy. Serial
terminal charging gives every $x\in[0,10]$ for every type assignment.
Since $q\ge7$, all mixed/new two-bus projected points are dominated by
the old two-bus interval. Therefore
\[
 D_Z=\min\{97,q+g(t)\},\qquad t=\max\{5,a\}.
\]
Every supported physical plan is a physical planner optimum. No projected
old plan can be rescued, and a dominated two-bus addition cannot have
zero regret: at equal cost it has an unsupported old counterpart; at
larger cost it has a strictly cheaper old counterpart at every price.
It suffices to test the new branch's unique physical minimizer $t$.

If $a\le5$, $t=5$, so $p=(5,5)$. All plans buy 30 units; comparison
with the old one-bus plan is simply $q\le7$, hence $q=7$.
If $a>5$, $t=a$ and $p_1-p_2=2a/5-2>0$. Within the new branch,
the endpoint $a$ minimizes the fixed-price objective. The best old
two-bus plan is its zero-early endpoint. Comparison against it gives
\[
 q\le14-a(p_1-p_2)=14+2a-\tfrac25a^2=A(a).
\]
Comparison against the old one-bus point gives
\[
 q\le7+(10-a)(p_1-p_2)=C(a).
\]
These comparisons include all physical alternatives by affine minimization
on every branch. They prove \eqref{eq:batteryregion}. For $a>5$,
$C(a)\ge7$ and $A(a)\ge7$ exactly when
$a\le(5+\sqrt{95})/2$ in the relevant domain; also $A(a)-C(a)=27-4a$.
At $q=7$, including $a\le5$, this proves \eqref{eq:batterythreshold}.

At fee 7 below the battery threshold, the hull's lower boundary joins
$(0,14)$ to $(a,7)$ and is then flat to 10. The hull objective on the
first edge is $14-7x/a+g(x)$ and is minimized at
$t_h=5+35/(2a)<a$. The physical optimum has early energy $a$;
its gap is $(a-t_h)^2/5$. At and above the threshold the physical
point $(\max\{5,a\},7)$ attains the hull minimum. Thus the boundary is
derived on the whole continuous battery domain, not inferred from a scan.

For $B=25$, the added endpoint is $(5,q)$. It is below the old lower
hull $14-0.7x$ exactly when $q<21/2$. In that regime the new lower
boundary joins $(0,14)$ to $(5,q)$ to $(10,7)$. The objective decreases
along the first edge. On the second it is minimized at
$x=5+(q-7)/2$ with operating cost $q-(q-7)^2/10$. Its value is
$85+q-(q-7)^2/20$. At $q=21/2$ this equals $7591/80$; above it the
whole new segment is dominated by the old hull. This proves the formulas.
For $q=8$, substitution gives $\CH_Z=1859/20$ and $D_Z=93$.
For $q=12$, both physical minima are 97. At the new load's price $(5,5)$,
the old one-bus bill is cheaper by 5. At the old load's price $(6,4)$,
the old zero-early two-bus bill remains 134; every new one-bus bill is
at least $q+130\ge137$. Thus the old plan's regret is 13, including
at the fee 12 tie. All quantities use fixed posted prices.
\end{proof}

\section{Replication proof}\label{app:replication}
\begin{proof}[Proof of Theorem~\ref{thm:replication} and
Corollary~\ref{cor:individual}]
For $m$ paired buses, the complete projected physical set has early energy
$E\in[10m,10n]$, late energy $30n-E$ and operating cost $14n-7m$.
The connector assignment in the text realizes every such early total:
paired buses take 10 each, and unpaired A buses distribute the remainder.
Convexification has lower operating boundary $14n-0.7E$. Substitution
in \eqref{eq:scaledF} gives the hull minimizer $E=27n/4$ and value
$7591n/80$.

For fixed integer $m\ge n/2$, supply is minimized at $E=10m$ and total
cost is
\[
 104n-27m+\frac{20m^2}{n}
 =\frac{7591n}{80}+\frac{20}{n}(m-27n/40)^2.
\]
For $m\le n/2$, the best early energy is $5n$ and branch cost is
$99n-7m$, decreasing with $m$. For odd $n$, moving from
$\lfloor n/2\rfloor$ to $\lceil n/2\rceil$ changes the minimized
cost by $-7+5/n<0$. For even $n$ the branch formulas meet at $n/2$.
Thus an integer nearest to $27n/40$ is optimal, including both at a tie,
and $|\delta|\le1/2$. This proves \eqref{eq:scaledgap}.

At that optimum the own-load marginal price is
\[
 p=(4+2m/n,\ 6-2m/n).
\]
The optimal $m$ has $m/n\ge1/2$, so a cheapest two-bus pair response
buys zero early (the whole interval ties when $m/n=1/2$).
The one-bus private objective minus that two-bus
objective is $40\delta/n$. If $\delta\ge0$, precisely the $m$
paired choices are disadvantaged; otherwise the $n-m$ unpaired choices
are. This proves \eqref{eq:scaledregret}. With one owner, all choices can
be replaced simultaneously. To compare separate ownership, first implement
the optimum using $m$ fixed-pair owners with one bus and $n-m$ with two.
Each one-bus owner buys 10 early; each two-bus owner buys zero early at the
planner optimum, giving the same total load and operating cost as the
aggregate optimum even if its literal matching crossed the new ownership
labels. Reserved connectors make each unilateral replacement feasible.
The identical projected opportunity sets give the same complete-fleet
response value, while each owner's gain is at most $20/n$. Summing the
disadvantaged owners' gains yields the whole-owner expression.

For $n=40k+1$, the optimum is $m=27k+1$, so $\delta=13/40$;
substitution yields the displayed subsequence. If $40$ divides $n$,
$\delta=0$. At $n=20$, the tied $m=13,14$ have $\delta=-1/2,1/2$,
giving regrets 7 and 14 respectively. All ties concern physical plans;
no fractional bus schedule is being implemented.
\end{proof}

\bibliographystyle{plainnat}
\bibliography{references}
\end{document}
